\section{Poker Head-to-Heads and the Value-Agreement Certificate}
\label{app:libratus}

The HUNL subgames and the Libratus endgames are built by our own code, so putting an
external solver on them requires handing it our tree in a format it can load.
This appendix describes how, what the resulting comparisons say, and the certificate
that establishes the three backends are solving the same game.

\paragraph{Exporting our trees.}
LiteEFG's \texttt{FileEnv} reads a generic extensive-form text format --- node lines for
terminal payoffs, chance edges with conditional probabilities, and player action lists,
plus infoset membership lines --- and its OpenSpiel bridge is one of several possible writers for that format. We emit it from our intermediate representation in a
single streaming pass with the root written first, so peak memory during export is set by the
infoset membership index, and the serialized text streams to disk. Kim (2026) reaches the same
file through a black-box game adapter, which is necessary because it runs in a separate
interpreter that must not import our package. Both backends load one checksummed
artifact per game, produced once and cached.

One detail of the adapter is worth recording because getting it wrong is silent. Kim
(2026) builds each infoset's action set by unioning the labels the adapter reports at
that infoset's member nodes, so a label must name an \emph{action}. Child
node names are globally unique and would therefore split each infoset into as many
action sets as it has members: the tree keeps its shape and every count still looks
right, while the solver optimizes over a strictly finer strategy space than the file
describes. On Kuhn poker that raises the sequence count from 13 to 25 per player and
improves apparent exploitability by more than two orders of magnitude, a number
comparable to nothing. Labelling actions by position among a node's children fixes it,
and is sound because the exporter validates that every member node of an infoset lists
its actions in the infoset's order.

\paragraph{Structural equivalence.}
Two checks confirm the export preserves the game. Exported Kuhn poker solved through the
adapter reproduces the exploitability of the same game loaded natively from OpenSpiel
\emph{bitwise}, both giving $0.0020750612020492554$ at $13$ sequences per player, and exported
Leduc poker gives $1093$ sequences per player either way. On the trees that matter here
the sequence counts Kim (2026) derives from the file agree exactly with what GPU-CFR's
compiler holds: the river subgame yields $701$ row and $701$ column sequences on
both sides, and Libratus endgame~4 yields $19{,}741$ and $27{,}497$, reproduced by
the baseline to the unit on a tree of $36{,}098{,}785$ nodes.

\paragraph{The certificate that scales: comparing on the value.}
The \emph{value} is comparable across backends whatever each one runs, because a
zero-sum tree has exactly one. For a
profile $(x,y)$ with per-player best-response gains $\delta_0, \delta_1$ and
$\mathrm{NashConv} = \delta_0 + \delta_1$, the usual sandwich
$\min_{y'} u_0(x,y') \le v^\ast \le \max_{x'} u_0(x',y)$ gives
$|u_0(x,y) - v^\ast| \le \max(\delta_0,\delta_1) \le \mathrm{NashConv}$, so for any two
backends
%
\begin{equation}
    |v_i - v_j| \;\le\; \mathrm{NashConv}_i + \mathrm{NashConv}_j .
    \label{eq:valuecert}
\end{equation}
%
\cref{eq:valuecert} is an exact identity: if it fails, at least one backend is
not solving the tree the other one is. It is also the check that scales, since it needs
only two scalars per backend and applies unchanged to trees far too large to compare
structurally. It is what caught the labelling bug above, and it is reported for every
poker head-to-head below.

\paragraph{Gaps agree once the variant is matched.}
The certificate constrains values, so on its own it leaves open whether the backends
merely solve the same game or also converge alike. Exploitability settles that, but only
when the algorithm is held fixed: LiteEFG's CFR$^{+}$ preset uses alternating updates
with linear averaging, and reading its gap against a simultaneous-update run compares
two algorithms. Matching the variant removes the
confound. Rerunning GPU-CFR on the river subgame in LiteEFG's own configuration
--- CFR$^{+}$, alternating, linear averaging, float64 --- gives exploitability
$1.27 \times 10^{-5}$ against LiteEFG's $1.25 \times 10^{-5}$ at $T = 8{,}000$, a
$1.6\%$ difference, with game values $1.77774849$ and $1.77774846$ agreeing to
$3 \times 10^{-8}$. Two independently written deterministic solvers reaching eight
significant figures of agreement on a $137{,}415$-state tree is a sharper statement than
\cref{eq:valuecert} alone can make, and it localizes the spread in the earlier tables to
the update order. The measurement is on the same
device and dtype for both, so it also fixes the price: an alternating iteration performs
one half-update per player and costs about twice a simultaneous one, and it still reaches
a smaller gap at matched wall-clock.

\paragraph{The river subgame, three ways.}
\cref{tab:riverthreeway} solves the $137{,}415$-state river subgame with all three
backends at two iteration counts. Every pair satisfies \cref{eq:valuecert} with wide
margin, using at most $24\%$ of the bound, and the agreement tightens with the gaps: an
eightfold increase in iterations shrinks the three-way value spread from
$1.02 \times 10^{-3}$ to $8.58 \times 10^{-5}$, a factor of $12$. The timings in the
table are median training-call wall-clock on the hardware of \cref{app:protocol};
quality metrics come from the fastest run. The
iteration counts are matched while the variants are not, so they measure what each
backend costs to run its own $1{,}000$ iterations, which is the quantity a
practitioner pays.

\input{figures/TABLE_T7_riverthreeway.tex}

\paragraph{Libratus endgame~4.}
The largest tree we put through all three backends is Libratus endgame~4 at full range:
$36{,}098{,}785$ nodes, exported as a $2.42$\,GB file. Its scale is what the export
mechanism buys and also what bounds it. GPU-CFR reads the endgame in $165.889$\,s,
compiles it and constructs the solver in $124.751$\,s, and solves it on one A100 at a
peak of $14.4$\,GiB allocated ($16.9$\,GiB reserved) in the benchmark process, which
also holds the float64 evaluation tree; the solver alone peaks at $8.9$\,GiB
(\cref{tab:lifecycle}). Kim (2026) parses the file in
$399.341$\,s and solves it on the same device; LiteEFG's
\texttt{FileEnv} loads it at $24.4$\,GiB resident and then iterates at
$13.856$\,s, so $200$ iterations is where all three meet. \cref{tab:sg4threeway}
(\cref{sec:results}) reports that point. The three values agree to $2.5$ chips on a
tree whose $200$-iteration gaps run to $77$ chips, using $1.4\%$ of what
\cref{eq:valuecert} allows, and the two GPU backends agree to $0.39$ chips at
$1{,}000$ iterations. The same $200$ iterations cost $3.180$\,s on GPU-CFR and
$2{,}771.135$\,s on LiteEFG: a factor of $870$ at matched iteration counts, each
backend running its own CFR$^{+}$, and the C++ implementation's advantage at small
scale (\cref{sec:matrix}) has inverted completely. On this tree the two GPU backends
cost within $14\%$ of each other: its $36$ million states carry $96{,}159$
Infosets, so our node-level dataflow and Kim's sequence-form products spend
$13.435$ and $11.820$\,ms per incremental iteration, and the $30$--$80\times$ gap of
\cref{tab:matrix} closes where the tree is almost all terminals. The GPU-CFR training times
report the median of two runs ($3.180$ and $3.188$\,s at $200$ iterations, $0.3\%$ apart),
with quality metrics from the fastest run; the baseline rows are single runs,
each costing $46$ minutes (LiteEFG) or a seven-minute parse (Kim). Each GPU-CFR
training call starts on a freshly constructed solver, so it includes the three eager
warmup iterations and graph capture; on this tree capture itself costs well under a
second (\cref{tab:lifecycle}), and the call is $197$ replays at the steady-state rate
plus that setup.

\paragraph{Endgame~3.}
We do not put the external baselines on endgame~3 ($194{,}556$ Infosets on
$93{,}168{,}610$ nodes): its
export would exceed $6$\,GB, and LiteEFG's load already peaks at $24.4$\,GiB on the
$2.6\times$ smaller endgame~4, so file loading would dominate the arm's time. The
head-to-head evidence stops at endgame~4's scale, where \cref{eq:valuecert} holds across
all three backends.
